\section{Introduction}
\label{sec:introduction}
Let $\T=\R/\Z$, let $\mu$ be normalized Lebesgue measure on $\T^3$,
and put $\ee(t)=\exp(2\pi i t)$. For an invertible measure-preserving
transformation $T$, we use the convention $U_Tg=g\circ T$.
The closed subspace of complex square-integrable functions of integral zero
is denoted by $L^2_0(\T^3,\mu)$.

On a probability space the constants form an invariant subspace, so the
probability-preserving form of Banach's simple Lebesgue-spectrum problem
asks for an invertible transformation whose Koopman operator on the
\emph{entire} orthogonal complement of the constants is unitarily
equivalent to multiplication by $w$ on $L^2(S^1,m)$, where $m$ is
normalized Haar measure. Equivalently, that restriction has a bilateral
orthonormal orbit which is complete. We realize this spectral model with
the standard smooth volume on a compact manifold.

\begin{theorem}\label{thm:main}
There exist a $C^\infty$ volume-preserving diffeomorphism
$T\colon\T^3\to\T^3$ and a real function $f\in L^2_0(\T^3,\mu)$ such that
\[
                  \{f\circ T^n:n\in\Z\}
\]
is an orthonormal basis of $L^2_0(\T^3,\mu)$.
Consequently $T$ is ergodic, and $U_T$ on this space is unitarily equivalent
to multiplication by the coordinate $w$ on $L^2(S^1,m)$.
\end{theorem}

The invariant measure in Theorem~\ref{thm:main} is the standard smooth
volume itself. The generator $f$ is only required to belong to $L^2$.
In particular, the assertion concerns the whole mean-zero space, including
both its spectral type and its multiplicity.
Section~\ref{sec:consequences} derives mixing, zero entropy, and vanishing
Lyapunov exponents. Mixing of every order uses, in addition, the separate
multiple-mixing theorem of \cite{MultipleMixing2026}.

The formulation above is distinct from the historical real-line question
recorded by Ulam and attributed to Banach \cite[\S 6, p.~76]{Ulam1960}.
Rokhlin asked for ergodic automorphisms with simple, or at least
finite-multiplicity, Lebesgue spectrum
\cite[\S 4, no.~7, p.~107]{Rokhlin1949}. Helson and Parry constructed
cocycles over aperiodic probability-preserving transformations whose
associated unitary operators have Lebesgue spectrum
\cite[p.~196, Theorem~(a)]{HelsonParry1978}. Mathew and Nadkarni constructed a
probability-preserving transformation with a Lebesgue \emph{component}
of multiplicity two \cite{MathewNadkarni1984}. Guenais connected the
Morse-cocycle question with flat polynomials and constructed a countable
abelian group action with simple spectrum of mixed discrete and Haar
type \cite{Guenais1999}. These results distinguish a Lebesgue component,
or a related group action, from one bilateral Lebesgue orbit spanning
the whole centered space.

Simple spectrum in smooth dynamics was already achieved by Fayad, who
constructed $C^\infty$ volume-preserving torus diffeomorphisms with
simple spectrum \cite{Fayad2001}; that result did not assert Lebesgue
spectral type.
In another setting, an infinite-measure conservative transformation
with simple Lebesgue spectrum is presented in
\cite[Theorem~2.2]{ElAbdalaoui2023}; the probability-space conclusion
does not follow from it. Theorem~\ref{thm:main} gives that conclusion
for the full mean-zero space of a smooth probability-preserving map.

Related existence results in continuous time have different spectral
consequences. Prikhod'ko constructs a finite-measure flow with simple
Lebesgue spectrum \cite[Theorem~3]{Prikhodko2020}. In the smooth category,
Fayad, Forni and Kanigowski obtain a real analytic conservative flow on
$\T^2$ with one singularity and countably infinite Lebesgue multiplicity
\cite[Theorem~1]{FFK2021}; Abdedou, Fayad and Kessi obtain real analytic
reparametrizations of minimal translation flows on $\T^5$ with the same
multiplicity \cite[Theorem~1.1]{AFK2023}. Taking a nonzero time map of a
simple Lebesgue flow yields countably infinite multiplicity: its spectral
model is multiplication by $\ee(t\xi)$ on $L^2(\R,\dd\xi)$, and the
substitution $\xi=(x+k)/|t|$, with $x\in[0,1)$ and $k\in\Z$, gives one
Lebesgue copy for each $k$. Thus these continuous-time results do not
supply Theorem~\ref{thm:main}.

Successive smooth changes of coordinates have a substantial history in
smooth dynamics, including the approximation-by-conjugacy constructions
of Anosov and Katok \cite{AnosovKatok1970}. We use explicit
volume-preserving passages between integrable twists rather than an
external conjugacy theorem. Their quantitative spectral-density
transport, and the finite Fourier signals built on it, are proved below.

\subsection*{The construction}
The approximating maps are integrable twists. In a volume-preserving
coordinate system on the torus, they have the form
\[
 S(\theta,u,a)=(\theta+\phi(u),u,a),
\]
where $\phi$ is a smooth real circle function that is strictly increasing
on an arc whose complement can be made arbitrarily short. A finite angular
Fourier expansion supported on that arc has an explicit absolutely
continuous spectral measure. We refer to its individual harmonics as
\emph{packets}, and keep a finite label, or \emph{color}, on disjoint groups
of packets. The colors carry the coefficients of an observable that we
want to approximate by translates of a single vector.

There are three steps. A slow passage in the coordinate $u$ replaces the
twist by an arbitrarily close twist whose angular coordinate is the old
$u$. Stationary phase locates the new angular indices and determines the
spectral density at a much finer scale. The spare coordinate $a$ supplies
a small adjustable variation in the speed (the \emph{jitter}). A large integer added to the
lift of $\phi$ separates the angular indices of different packets without
changing the old map. This separation is what allows later perturbations
to distinguish colors even when their temporal spectra overlap.

Next, a positive Fourier construction and a finite prime sieve produce a
small sinusoidal bias in each color density. Its coefficient depends on
the normalized angular index. A neutral passage between successive biased
passages makes residue classes carry equal shares of each color. The
bias strength can be chosen before the angular scale grows; the passage
parameter is then free to enforce any required smooth or spectral
accuracy.

Finally, many small biases are grouped into blocks. If color $r$ has
spectral density $p_r$ and target coefficient $b_r$, the center at a
temporal frequency $z$ is the density-weighted mean
\[
 p(z)=\sum_r p_r(z),\qquad
 \beta(z)=\frac{\sum_r b_rp_r(z)}{p(z)}\quad(p(z)>0).
\]
Centering each block at a smooth approximation to this current mean
makes the loss in $\int\sqrt p$ quadratic
in the block's strength budget, while the same block provides an estimator
of the color label. Repetition makes prediction arbitrarily accurate with
arbitrarily little loss in $\int\sqrt p$. This separation of the two
estimates is the central finite construction. When the vector's norm is
close to one, keeping $\int\sqrt p$ close to one makes whitening a small
change in physical $L^2$, while restoring the spectral density exactly
to one. A diagonal limit preserves the orthonormal orbit and makes its
span contain every member of a dense family of observables.

\subsection*{Order of the proof}
Section~\ref{sec:passages} proves the analytic passage rules, including the
comparison of jitter laws along a rapidly oscillating phase.
Section~\ref{sec:signals} constructs the positive laws that transmit the
small biases. Section~\ref{sec:blocks} proves the finite prediction result,
and Section~\ref{sec:completion} prepares arbitrary observables and takes
the smooth limit. Section~\ref{sec:consequences} derives the dynamical
consequences. Every passage acts on finite data. Its integer parameter
is chosen last, after the current coordinate chart, cutoffs, jitter,
and all desired error bounds have been fixed. No uniform bound on the
number of packets or on the derivatives of successive coordinate charts
is assumed.
